%% kapa: exam=1 round=25 year=2014 period=1 form=A subject=economics subject_ko=경제원론 seq=29 number=69
%% source: Q-Net 감정평가사 기출문제 (cst00309) artlSeq=5101207 file=extracted/1008/5101207/2014년 제25회 감정평가사 1차시험 기출문제/1교시(A형).hwp
%% conversion: kordoc 4.24.51 -> tools/gen1.py (auto); review=proofread-sonnet; flags=none
\begin{question}{69}
소비의 항상소득가설과 생애주기가설에 관한 설명으로 옳은 것을 모두 고른 것은?
\begin{framedtext}
ㄱ. 소비자들은 가능한 한 소비수준을 일정하게 유지하려는 성향이 있다. \par
ㄴ. 생애주기가설에 의하면 고령인구의 비율이 높아질수록 민간부문의 저축률이 하락할 것이다. \par
ㄷ. 프리드만(Friedman)의 항상소득가설에 의하면 높은 소득의 가계가 평균적으로 낮은 평균소비성향을 갖는다. \par
ㄹ. 케인즈(Keynes)는 항상소득가설을 이용하여 승수효과를 설명하였다. \par
\end{framedtext}
\begin{choices}
  \item ㄱ, ㄴ
  \item ㄱ, ㄹ
  \item ㄴ, ㄷ
  \item ㄱ, ㄴ, ㄷ
  \item ㄴ, ㄷ, ㄹ
\end{choices}
\end{question}
\answer{4}
